\clearpage
\fancyhead[L]{\small ELECTRODYNAMICS.}
\thispagestyle{empty}
\pdfbookmark[1]{Electrodynamics. Edition and source}{electrodynamics-edition}
\begin{center}\Large\textbf{On the Mathematical Theory of Electrodynamic Effects}\end{center}

This AI translation into English, with AI-produced typesetting, presents the title of the manuscript division and sections [1]--[21] of Carl Friedrich Gauss’s work in \textit{Werke}, vol.~V (Göttingen, 1867), pp.~601--630. Sections [1]--[3] concern current intensity in a network, the resistance and intensities in a particular circuit, and the action of two current elements on each other. Sections [4] and [5] concern the linking of two lines and the law of the galvanic current. Section [6] treats the force between current elements and the electromotive force. Sections [7] and [8] set out the law of induction and the laws for two closed conductors. Section [9] concerns induction by motion. Sections [10]--[12] give expressions for the effects of galvanic currents, treat polygons with straight sides, and relate electrical and magnetic quantities. Sections [13] and [14] present two formulations of the fundamental law for galvanic currents. Sections [15]--[17] concern the surface of a sphere, the electromotive force produced by a newly formed current, and the effect of a wound ring. Section [18] treats equilibrium and the action of a galvanic current on a body carrying positive and negative masses. Section [19] gives Gauss’s proof of Ampère’s fundamental theorem. Section [20] is Gauss’s letter to Wilhelm Weber of 19 March 1845. Section [21] lists relations among physical quantities. The diagrams are drawn in editable \LaTeX{}.

\medskip\noindent
\href{https://gdz.sub.uni-goettingen.de/id/PPN236006339}{Collected German edition (1867)}

The 1833 note on linking also appears in English translation in Renzo L. Ricca and Bernardo Nipoti, \textit{Gauss’ Linking Number Revisited}, Fig.~1(b), p.~1327, and in Martin Mathieu’s translation in Charles Nash, \textit{Topology and Physics---a Historical Essay}, §~1, pp.~3--4, note~3. The text and formula here follow the collected German edition.

\medskip\noindent
\href{https://webhomes.maths.ed.ac.uk/~v1ranick/papers/ricca.pdf#page=3}{Ricca and Nipoti (2011)}\quad
\href{https://arxiv.org/abs/hep-th/9709135v3}{Mathieu’s translation in Nash (1997)}
\par\medskip\noindent Clemens Schaefer gives the expanded wording of the discovery time in the dated note to section [7] in \href{https://www.nature.com/articles/128339a0}{\textit{Gauss’s Investigations on Electrodynamics}} (1931), p.~341, note~2.
\par\medskip\noindent The \href{https://archive.org/details/Werkecarlf05gausrich/page/615/mode/1up}{second printing of \textit{Werke}, vol.~V (Göttingen, 1877), p.~615}, gives the symbol key in section [10] and the logarithmic expressions in section [11].
\par\medskip\noindent A.~K.~T. Assis translated section [13] in \href{https://www.ifi.unicamp.br/~assis/Weber-in-English-Vol-4.pdf#page=14}{\textit{Wilhelm Weber’s Main Works on Electrodynamics Translated into English}, vol.~IV (Apeiron, 2021), chapter~3, p.~11}. James Clerk Maxwell quotes the opening sentence in English in \href{https://en.wikisource.org/wiki/Page:A_Treatise_on_Electricity_and_Magnetism_-_Volume_2.djvu/461}{\textit{A Treatise on Electricity and Magnetism}, vol.~II (1873), §~851, p.~429}.
\par\medskip\Needspace{4\baselineskip}\noindent An English translation of section [15] by user2554 is available in the \href{https://hsm.stackexchange.com/questions/7503/several-questions-about-gausss-contributions-to-electromagnetism}{15 August 2025 revision of a History of Science and Mathematics Stack Exchange question}. The diagrams of sections [16] and [17] also appear in the \href{https://archive.org/details/Werkecarlf05gausrich/page/621/mode/1up}{second printing of \textit{Werke}, vol.~V (1877), p.~621}. Maxwell cites the result of section [17] in \href{https://en.wikisource.org/wiki/Page:A_Treatise_on_Electricity_and_Magnetism_-_Volume_2.djvu/344}{\textit{A Treatise on Electricity and Magnetism}, vol.~II (1873), §~706, p.~312}.
\par\medskip\noindent Susan P. Johnson’s English translation of the letter in section [20] appears in \href{https://21sci-tech.com/articles/Atomic_Science.pdf#page=22}{\textit{Text of the Gauss-Weber 1845 Correspondence}}, edited by Laurence Hecht, \textit{21st Century Science \& Technology} 9(3), Fall 1996, pp.~42--43. Hecht wrote the accompanying article. The editorial note on p.~41 credits Karl Krause and Alexander Hartmann for transcribing the three adjacent letters from Weber to Gauss.
\par\smallskip\noindent In Gauss’s letter, \textit{in verbis facilior} means “more accommodating about wording,” and \textit{Nil actum reputans si quid superesset agendum} means “counting nothing as done while anything remained to be done.”
\par\smallskip\noindent Section [21] also appears in the \href{https://archive.org/download/Werkecarlf05gausrich/page/n641}{second printing of \textit{Werke}, vol.~V (1877), p.~630}.
\clearpage
